% Part: lambda-calculus
% Chapter: lambda-definability
% Section: truth-values

\documentclass[../../../include/open-logic-section]{subfiles}

\begin{document}

\olfileid{lam}{ldf}{tvr}
\olsection{સત્યમૂલ્યો અને સંબંધો}

શુદ્ધ લૅમ્બડા કલનમાં સત્યમૂલ્યો આ રીતે
સંકેતિત કરી શકાય:
\begin{align*}
  \fn{true} & \ident \lambd[x][\lambd[y][x]]\\
  \fn{false} & \ident \lambd[x][\lambd[y][y]]
\end{align*}

સત્યમૂલ્યોનું નિરૂપણ \emph{પસંદગી-વિધેયો} તરીકે
થાય છે, એટલે બે આર્ગ્યુમેન્ટ લઈને તેમાંથી એક
પાછું આપતાં વિધેયો તરીકે. સત્યમૂલ્ય
$\fn{true}$ પહેલો આર્ગ્યુમેન્ટ પસંદ કરે છે અને
$\fn{false}$ બીજો. દાખલા તરીકે,
$\fn{true}\, M N$ હંમેશા~$M$માં લઘુકૃત થાય છે,
જ્યારે $\fn{false}\, M N$ હંમેશા~$N$માં થાય છે.

\begin{defn}
સંબંધ $R \subseteq \Nat^k$ને !!{lambda definable}
કહીએ, જો એવું પદ~$R$ હોય કે
\begin{align*}
  R\, \num{n_1} \dots \num{n_k} & \bred \fn{true}
  \intertext{જ્યારે $R(n_1, \dots, n_k)$ સાચો હોય, અને}
  R\, \num{n_1} \dots \num{n_k} & \bred \fn{false}
\end{align*}
બાકી કિસ્સામાં.
\sourcecorrection{OLLAM-053}{સ્થિર અંગ્રેજી મૂળમાં
સંબંધનું ક્ષેત્ર~$\Nat^n$ છે, પરંતુ તેની દલીલો
$n_1,\dots,n_k$ અને અનુપ્રયોગમાં પણ~$k$
ચર્ચ અંકો છે. અરીટી એકરૂપ રાખવા અહીં
$R\subseteq\Nat^k$ લખ્યું છે.}
\end{defn}

દાખલા તરીકે, જે~$0$ માટે અને ફક્ત~$0$ માટે જ સાચો છે તે સંબંધ
$\fn{IsZero} = \{0\}$ આ પદ વડે
!!{lambda definable} છે:
\[
  \fn{IsZero} \ident \lambd[n][n (\lambd[x][\fn{false}])\, \fn{true}].
\]
તે કેવી રીતે કામ કરે છે? ચર્ચ અંકો પુનરાવર્તકો
તરીકે વ્યાખ્યાયિત છે: બીજો આર્ગ્યુમેન્ટ લઈને
તેના પર પહેલા આર્ગ્યુમેન્ટને~$n$ વાર લાગુ કરે છે.
અહીં શરૂઆતનું મૂલ્ય~$\fn{true}$ છે, અને
પુનરાવર્તનના દરેક પગલે અગાઉના પગલાનું પરિણામ
ગમે તે હોય, $\fn{false}$ મળે છે. આ પગલું
શરૂઆતના મૂલ્ય પર $n$ વાર લાગુ થાય છે. એટલે
પગલું એકેય વાર ન લાગ્યું હોય, એટલે કે $n = 0$
હોય, ત્યારે અને માત્ર ત્યારે જ પરિણામ
$\fn{true}$ મળે છે.

સત્યમૂલ્યોના આ નિરૂપણ પરથી બીજા સત્ય-વિધેયો
પણ વ્યાખ્યાયિત કરી શકાય. અહીં નિષેધ અને
સંયોજનનાં બે નિરૂપણ છે:
\begin{align*}
  \fn{Not} & \ident \lambd[x][x\, \fn{false}\, \fn{true}]\\
  \fn{And} & \ident \lambd[x][\lambd[y][xy \,\fn{false}]]
\end{align*}
$\fn{Not}$ એક આર્ગ્યુમેન્ટ લે છે અને તે
$\fn{false}$ હોય તો $\fn{true}$, તથા
$\fn{true}$ હોય તો $\fn{false}$ આપે છે.
$\fn{And}$ બે સત્યમૂલ્યો લે છે અને બંને
$\fn{true}$ હોય ત્યારે અને માત્ર ત્યારે જ
$\fn{true}$ આપવું જોઈએ. ઉપર જણાવ્યા પ્રમાણે
સત્યમૂલ્યો પસંદગી-વિધેયો છે: $x$ સત્યમૂલ્ય હોય
અને તેને બે આર્ગ્યુમેન્ટ લાગુ કરીએ તો~$x$
$\fn{true}$ હોય ત્યારે પહેલો, અન્યથા બીજો
આર્ગ્યુમેન્ટ મળે. હવે $\fn{And}$ પોતાના
આર્ગ્યુમેન્ટો $x$ અને~$y$ લઈને પહેલા
આર્ગ્યુમેન્ટ~$x$ને $y$ અને $\fn{false}$
સોંપે છે. $x$ સત્યમૂલ્ય હોય એમ માનીએ તો
$x$ એ $\fn{true}$ હોય ત્યારે પરિણામ~$y$,
અને $x$ એ $\fn{false}$ હોય ત્યારે પરિણામ
$\fn{false}$ મળે છે. આ જ જોઈએ છે.

અહીં આપણે માની લીધું છે કે $\fn{And}$ને
આર્ગ્યુમેન્ટ તરીકે માત્ર સત્યમૂલ્યો જ અપાય છે.
તેને બીજાં પદો આપીએ તો પરિણામ, એટલે સામાન્ય
સ્વરૂપ અસ્તિત્વમાં હોય તો તે, સત્યમૂલ્ય ન પણ હોય.

\begin{prob}
  સત્યમૂલ્યોના $\lambd$-પદ તરીકેના આ સંકેતન
  વડે સમાવેશી તથા અપવર્જી વિયોજનના સત્ય-વિધેયો
  દર્શાવતાં $\fn{Or}$ અને $\fn{Xor}$ વ્યાખ્યાયિત કરો.
\end{prob}

\end{document}
